\documentclass{report}
\usepackage{amsmath,amssymb}
\setlength{\parindent}{0mm}
\setlength{\parskip}{1em}
\begin{document}
\begin{center}
\rule{6in}{1pt} \
{\large Kees Vuik \\
{\bf Physics-based pre-conditioners for large-scale subsurface flow simulation}}

Delft University of Technology \\ Delft Institute of Applied Mathematics \\ Mekelweg 4 \\ 2628 CD Delft \\ The Netherlands
\\
{\tt c.vuik@tudelft.nl}\\
Gabriela Diaz\\
Jan Dirk Jansen \end{center}

\nonumber
\noindent
Physics-based pre-conditioners for large-scale subsurface flow
simulation
\\[2ex]
C. Vuik, G. Diaz and J.D. Jansen\\
Delft University of Technology\\
Department of Geoscience and Engineering
and
Delft Institute of Applied Mathematics\\
Mekelweg 4,
2628 CD Delft,
The Netherlands \\
{\bf e-mail: c.vuik@tudelft.nl}\\
[2ex]
In this paper we consider new physics-based pre-conditioning techniques
for solving large, sparse linear systems of equations with strongly
varying coefficients and localised source terms.
Such systems occur during the simulation of multi-phase flow through
strongly heterogeneous porous media as occur in e.g. oil and gas
production, CO$_2$ storage or (contaminated) ground water flow.
The underlying partial differential equations are typically of two
different types: a weakly nonlinear near-elliptic parabolic equation
with time-varying parameters to describe the diffusive pressure
behaviour, and one or more strongly nonlinear nearly-hyperbolic parabolic
equations to describe the mostly convective behaviour of the interfaces
between the phases.
The spatially discretized coupled nonlinear equations are typically
solved with a low-order time discretization and Newton-Raphson
iteration.
\\[2ex]
The state of the art approach to solve
these systems is to apply a two-stage preconditioning in which first the
pressure-related part of the matrix is isolated (with the aid of Schur
decomposition) and solved, typically with an Algebraic Multi Grid (AMG)
solver. Thereafter the full system of equations is solved with e.g. a
Conjugate Gradient solver.
We consider the first step of this two-step procedure, i.e. solving of
the pressure equation. A particular aspect of porous media flow is the
relatively small number of wells which are usually positioned quite far
apart. This results in pressure solutions that are dominated by the
near-well pressure behaviour in addition to the geological spatial
features.
\\[2ex]
In [1] it is shown that an effective alternative to the AMG solver for
the pressure equation can be provided by a
reduced-order model obtained by applying proper orthogonal decomposition
(POD) to a number of pre-computed solutions.
In case of optimisation or
parameter estimation exercises which require tens to hundreds of
'reservoir simulations' of nearly identical models with nearly identical
right-hand sides (source terms), the CPU time to compute the
reduced-order model is negligible.
\\[2ex]
In this paper we analyse and further develop
the preconditioners given in [1]. We investigate alternatives using
different varieties of reduced-order modelling. Furthermore we explore
the connection between POD-based preconditioning and deflation methods.
One of the difficulties for deflation methods is to find the right
deflation vectors for general problems. A start has been made to find
these vectors in an automated way, however until now this remains an
open question. The combination of deflation with the POD methods looks
very promising in this respect.
\\[2ex]
Some numerical experiments are given to illustrate the theory. We start
with a simple layered problem and also give results for the SPE10
benchmark [2].
\section*{References}
[1]
Astrid, P., Papaioannou, G.V., Vink, J.C. and Jansen, J.D., 2011:
Pressure preconditioning using proper orthogonal decomposition. Paper
SPE 141922 presented at the SPE Reservoir Simulation Symposium, The
Woodlands, USA, 21-23 February
\\[2ex]
[2] http://www.spe.org/web/csp/datasets/set02.htm


\end{document}
